\section{System Architecture}
\label{sec:architecture}

The engine has one job: run a fixed catalogue of econometric methods, return each
result as structured data, and attach enough provenance that the number can be
banded against an expectation and reproduced later. Everything in this section
follows from a single decision -- that a caller may name a method and supply
parameters, but may never supply code. That decision is what makes the engine safe
to expose publicly, and it is what lets the same method run identically in a cloud
container, on the laboratory workstation, and in the test harness.

We organise the system into three planes, each chosen for a different operational
regime. The \emph{kernel} answers synchronous requests on scale-to-zero Cloud Run,
which suits bursty, low-volume research traffic at near-zero idle cost. The
\emph{tower} carries the heavy asynchronous estimators that a request-scoped cloud
service cannot, on an always-on workstation with governed cores and a GPU. The
\emph{surfaces} are static pages that render only what the evaluation suite has
verified, so the read side computes nothing and can never present an unchecked
figure. Figure~\ref{fig:three-planes} sketches the topology.

\begin{figure}[t]
\centering
\begin{tikzpicture}[font=\small,
  box/.style={draw=cegrey, thick, rounded corners=2pt, align=center, inner sep=5pt},
  lbl/.style={font=\footnotesize\itshape, text=cegrey},
  flow/.style={-{Stealth[length=2.4mm]}, thick, draw=cegrey}]
  \node[box, fill=ceblue!12, text width=10.6cm] (read) at (0,2.5)
    {\textbf{Surfaces} \,\textnormal{(static GitHub Pages)}\\[1pt]
     \footnotesize research dashboard $\cdot$ \file{evals.html} $\cdot$ \file{reproduce.html} $\cdot$ \file{claims.html} $\cdot$ gallery\\
     \footnotesize render \file{evals.json} -- no compute on the read side};
  \node[box, fill=cegreen!14, text width=4.6cm] (kernel) at (-3.2,0)
    {\textbf{Kernel}\\[1pt]
     \footnotesize \file{compute-server.mjs}\\
     \footnotesize Cloud Run \textit{shssm-compute}\\
     \footnotesize sync methods, \route{/api/chat},\\ \route{/api/claims}, \route{/health}};
  \node[box, fill=ceorange!14, text width=4.9cm] (tower) at (3.2,0)
    {\textbf{Tower}\\[1pt]
     \footnotesize \file{job-server.mjs} :3030\\
     \footnotesize always-on workstation, \code{MAX\_CONCURRENT}=2\\
     \footnotesize 4 heavy async methods\\ \footnotesize governed cores + GPU};
  \node[lbl, below=2pt of kernel] {scale-to-zero};
  \node[lbl, below=2pt of tower] {always-on};
  \draw[flow, dashed] (read.south) -- (kernel.north)
    node[midway, sloped, above, font=\footnotesize] {fetch (read-only)};
  \draw[flow, {Stealth[length=2.4mm]}-{Stealth[length=2.4mm]}] (kernel.east) -- (tower.west)
    node[midway, above, font=\footnotesize] {submit / poll};
\end{tikzpicture}
\caption{The three planes. The read side fetches but never computes; the kernel
serves synchronous methods and submits heavy work to the tower.}
\label{fig:three-planes}
\end{figure}

\subsection{The kernel: synchronous methods on scale-to-zero Cloud Run}
\label{sec:kernel}

The kernel is \file{server/compute-server.mjs}, a zero-dependency Node service that
needs only \code{node:http}, \code{node:https}, \code{node:child\_process} and
\code{node:fs}. It is deployed to Google Cloud Run as the service
\code{shssm-compute} in region \code{asia-south1}; the live release is revision
\code{shssm-compute-00036-llf}, re-confirmed on 2026-06-15 serving a 26-method
registry. A liveness probe reports the running configuration directly:
\route{GET /health} returns
\code{\{ok:true, methods:26, revision:"shssm-compute-00036-llf",
sandbox:"timeout", timeout\_s:90\}}. The same code runs locally on port
\code{3200} for development.

Cloud Run scales the kernel to zero between requests. For a research service whose
traffic is occasional and bursty -- a workbench query here, a chatbot turn there --
this is the right trade: we pay for compute only while a request is in flight, and
a cold start is acceptable latency for an interactive analysis. The cost of that
choice is the request budget. Cloud Run bounds a request at 300\,s and may throttle
CPU outside an active request, so any estimator whose honest runtime is measured in
minutes cannot live on the synchronous path. Those estimators are the tower's
(Section~\ref{sec:tower}); the kernel carries the rest, which return in well under
the timeout.

The kernel's public surface is a small, rate-limited JSON API. CORS is open so the
static read-side pages can fetch it, but those pages only ever read; nothing on the
read side mutates engine state. Table~\ref{tab:kernel-surface} reproduces the
routes. The two model-backed routes, \route{/api/chat} and \route{/api/research},
share a per-instance daily budget and, over cap, return \code{503 daily\_capacity}
with a reset hint rather than a fabricated answer; that budget and its cost ceiling
are treated in Section~\ref{sec:deployment-ops}.

\begin{table}[t]
\centering
\small
\caption{The kernel HTTP surface (\file{server/compute-server.mjs}). Read-side
pages only ever consume these routes; the per-minute and per-day caps are the
abuse and credit guards of Section~\ref{sec:deployment-ops}.}
\label{tab:kernel-surface}
\begin{tabular}{@{}p{3.5cm}l p{7.4cm} l@{}}
\toprule
Route & Method & Purpose & Cap \\
\midrule
\route{/api/compute/run}     & POST & run one registry method on validated params      & 20/min \\
\route{/api/compute/catalog} & GET  & method + dataset registry the workbench binds to  & -- \\
\route{/api/chat}            & POST & AI analyst (Gemini function-calling over registry) & 10/min, 50/day \\
\route{/api/research}        & POST & deep-research assistant (agentic + grounded)       & 5/min, 20/day \\
\route{/api/situate}         & POST & place a result against the verified record         & 30/min \\
\route{/api/claims}          & GET  & epistemic ledger (established / contested / superseded) & -- \\
\route{/api/sri/current}     & GET  & systemic-risk live-feed: latest value              & -- \\
\route{/api/sri/history}     & GET  & systemic-risk live-feed: time series               & -- \\
\route{/api/sri/network}     & GET  & systemic-risk live-feed: directed edges            & -- \\
\route{/api/sri/cron-tick}   & POST & the daily SRI append (Cloud Scheduler only)        & -- \\
\route{/api/event}           & POST & read-side telemetry (aggregate counts, no PII)     & 60/min \\
\route{/health}              & GET  & liveness: methods, revision, sandbox               & -- \\
\route{/metrics}             & GET  & request / cache / rate-limit counters             & -- \\
\bottomrule
\end{tabular}
\end{table}

\subsection{The tower: heavy asynchronous estimators on an always-on workstation}
\label{sec:tower}

Four methods cannot live on the synchronous path, because their honest runtimes
exceed what a scale-to-zero, request-scoped service can carry: \method{ksg\_te}
(exact KSG/Frenzel--Pompe transfer entropy with IAAFT surrogates across all
directed pairs), \method{ksg\_robustness} (its nuisance-grid sweep),
\method{namh\_pipeline} (the end-to-end NAMH estimation, roughly 24 minutes at
$B=200$), and \method{soch\_robustness} (a bootstrap symmetry-test grid). These run
on the tower, \file{server/job-server.mjs}, an always-on service on port
\code{3030} on the group's 22-core workstation, which also carries the GPU used to
offload the transfer-entropy search (treated in the heterogeneous-compute section).

The tower inverts the kernel's request model. \route{POST /api/jobs/submit}
validates the method and returns a \code{job\_<date>\_<hash>} id immediately with
\code{202 Accepted}; the work then runs asynchronously at a small fixed concurrency
(\code{MAX\_CONCURRENT = 2}). A caller polls \route{GET /api/jobs/\{id\}} for the
record, or subscribes to \route{GET /api/jobs/\{id\}/stream}, a Server-Sent-Events
channel that emits \code{progress} events and terminates on \code{done} or
\code{error}. The progress events are exactly the \code{CE\_PROGRESS}-gated lines of
Section~\ref{sec:json-contract}: the tower spawns the same runner with
\code{CE\_PROGRESS=1}, parses each newline-delimited progress object to update the
job, and keeps the final unterminated object as the result.

Job state persists to disk. Every job record is written to \file{.jobs/} as a JSON
file on each transition, so the tower survives a restart: on boot it re-reads
\file{.jobs/}, and any job left \code{running} or \code{queued} is re-queued rather
than lost. Crucially, the tower is also where core governance lives -- it sets a
hard per-job core ceiling and pins the numerical libraries to a single thread each
before spawning R, so two concurrent jobs cannot oversubscribe the machine. That
mechanism, and the 22-core saturation it fixed, are the subject of the
heterogeneous-compute section; here it suffices that the budget is set by the plane
that owns concurrency.

\subsection{The surfaces: render only what is verified}
\label{sec:surfaces}

The third plane is the read side: static GitHub Pages in the sibling
\code{econstellar/} checkout -- the research dashboard, the evaluation and
reproduction pages, the claims panel, the gallery. Nothing on the read side
computes. Each page fetches and renders \file{evals.json}, the single
machine-written artefact of one genuine run of the public evaluation suite, and the
kernel's read-only routes (\route{/api/claims}, the SRI feed). A figure appears on a
surface only because the suite produced it under a pre-registered, failable band; a
regression turns the dashboard red rather than passing silently.

\begin{principle}[Render only what is verified]
\label{prin:render-verified}
The read surfaces display \file{evals.json} and the kernel's read routes; they
execute no analysis of their own. A number reaches a reader only after the
evaluation suite has produced it and checked it against a band fixed before the run.
The read side cannot present an unchecked figure because it cannot compute one.
\end{principle}

This is the architectural counterpart of the catalogue discipline. The kernel will
not run a method that is not registered; the surfaces will not show a result that
the suite has not verified. The separation also keeps the planes honest about cost
and blast radius: the read side is free static hosting that can fail closed, the
kernel pays only while computing, and the heavy machine is reached only through a
job submission that the kernel itself gates.

\subsection{The request lifecycle}
\label{sec:lifecycle}

Every synchronous analysis follows one path, and the path is the same whether the
caller is the workbench, the chatbot, or the test harness. We trace it concretely
for \route{POST /api/compute/run}.

\begin{enumerate}[leftmargin=2em,itemsep=2pt]
\item \textbf{Receive.} The body is read with a 64\,KB cap; an oversized body is
  rejected with \code{413} before parsing, so no request can fan out into unbounded
  memory. A malformed body returns \code{400 bad JSON body}.
\item \textbf{Validate.} \code{validate(method, params)} looks the method up in the
  fixed registry. An unknown name returns \code{400 unknown method}. Each declared
  parameter is then coerced and checked against its typed schema (Section~%
  \ref{sec:registry-guard}); a market name not in the panel's whitelist is rejected
  as \code{unknown series}. What reaches R is a \emph{clean} parameter object, never
  the raw request.
\item \textbf{Route by kind.} If the method is flagged \code{long\_running} it is
  refused here with \code{\{error:"...background job...", async:true\}} and directed
  to the tower; the synchronous endpoint never spawns a heavy estimator.
\item \textbf{Cache.} A 5-minute response cache, keyed on method and cleaned
  parameters, short-circuits repeated identical analyses. Errors are never cached.
\item \textbf{Spawn.} On a cache miss the kernel builds the runner argument vector
  and spawns the registered script under the sandbox (Section~\ref{sec:sandbox}):
  \code{Rscript r/<script>.R '<params-json>'}, wrapped in \code{timeout}. The R
  helper reads exactly the whitelisted dataset and computes.
\item \textbf{Collect.} The child's \code{stdout} is accumulated and parsed as a
  single JSON object (Section~\ref{sec:json-contract}). A \code{124} exit is
  surfaced as an explicit \code{timeout after Ns} error, never a silent hang; a
  non-JSON stdout is reported as such rather than passed through.
\item \textbf{Respond.} The result is wrapped with a provenance stamp -- method
  version, engine revision, cleaned parameters, panel id and byte-hash, timestamp,
  and a permalink that re-runs the exact analysis on load -- and returned.
\end{enumerate}

\begin{figure}[t]
\centering
\resizebox{\textwidth}{!}{%
\begin{tikzpicture}[font=\footnotesize, node distance=6mm,
  st/.style={draw=cegrey, thick, rounded corners=2pt, fill=celight, align=center,
    inner sep=4pt, minimum height=9mm, text width=1.75cm},
  ar/.style={-{Stealth[length=2mm]}, thick, draw=cegrey}]
  \node[st] (req) {HTTP\\request};
  \node[st, right=of req] (guard) {size\\guard};
  \node[st, right=of guard] (val) {validate vs.\\registry};
  \node[st, right=of val] (disp) {sync /\\async?};
  \node[st, right=of disp] (spawn) {sandbox\\$+$ \file{Rscript}};
  \node[st, right=of spawn] (json) {single-JSON\\$+$ provenance};
  \node[st, right=of json] (resp) {response};
  \foreach \a/\b in {req/guard, guard/val, val/disp, disp/spawn, spawn/json, json/resp}
    \draw[ar] (\a) -- (\b);
  \node[st, below=8mm of disp, fill=ceorange!14, text width=1.9cm] (tower) {tower job (async)};
  \draw[ar] (disp) -- (tower);
  \node[font=\scriptsize, text=ceorange, below=0.5mm of val.south] {400 unknown};
  \node[font=\scriptsize, text=ceorange, below=0.5mm of spawn.south] {124 $=$ timeout};
\end{tikzpicture}}
\caption{The synchronous request lifecycle. Validation against the registry
precedes any spawn; a method flagged long-running is refused synchronous service
and submitted to the tower; a wall-clock timeout surfaces as an explicit error,
never a silent hang.}
\label{fig:lifecycle}
\end{figure}

Listing~\ref{lst:lifecycle} states the same path as pseudocode, to make the
ordering of the guards unambiguous: validation precedes any spawn, and the registry
lookup is the only thing that turns a request name into an executable.

\begin{lstlisting}[language=JavaScript,caption={The synchronous request path, in
brief. Validation and registry lookup precede every spawn; the runner receives a
cleaned parameter object, never caller-supplied code.},label={lst:lifecycle}]
// POST /api/compute/run  (server/compute-server.mjs)
if (bodyBytes > 64*1024)        return send(413);          // size guard
let payload = JSON.parse(body);                            // 400 on bad JSON
let { method, clean } = validate(payload.method,           // 400 unknown method
                                 payload.params);          // typed-param coercion
if (method.long_running)        return send(400, async);   // -> tower, never here
let hit = cacheGet(method, clean);                         // 5-min TTL
let result = hit
  || await runSandboxed(method, buildArgsJson(method, clean));
return send(200, { ...result,
                   provenance: stamp(method, clean) });    // version+revision+hash
\end{lstlisting}

\subsection{The registry as the primary safety property}
\label{sec:registry-guard}

The central design fact is that the registry is parameterised only. A method entry
is a fixed runner script plus a typed parameter schema -- for example, the
stationarity method binds \code{r/adf.R} with a \code{series} parameter that must be
exactly one market drawn from the panel's whitelist. The runner scripts in
\file{r/} are part of the repository and are never user-supplied. A caller therefore
chooses \emph{which} registered analysis to run and supplies typed inputs to it, but
has no channel through which to inject code, a file path, or a shell fragment.

\begin{invariant}[Parameterised-only execution]
\label{inv:no-rce}
The only executables the engine ever runs are the registered runner scripts shipped
in the repository. A request supplies a method name from a closed registry and a
parameter object validated against that method's typed schema. No request input is
interpreted as code, a path, or a command. There is therefore no
arbitrary-code-execution surface, independent of any sandbox.
\end{invariant}

This is the load-bearing claim, and it is worth being precise about why. The
validation step (Section~\ref{sec:lifecycle}) coerces each parameter to its declared
type and rejects anything outside its domain: series names are checked against the
dataset whitelist, integers and numbers must parse finitely, enumerations must match
a fixed value set, and counts must fall in the method's declared arity. A string that
is not a registered method name does not select a default; it is refused. A probe
such as \code{"DROP\_TABLE; rm -rf /"} as a method name returns \code{400 unknown
method} and reaches no shell. The safety of the public endpoint rests here, on
Invariant~\ref{inv:no-rce}, not on the sandbox.

The sandbox is defence in depth, and we are careful not to overstate it
(Section~\ref{sec:sandbox}). Stating the property in the order it actually holds --
registry first, sandbox second -- matters because the two deployments run different
sandbox modes, and the safety argument must not depend on which one is live.

\subsection{The sandbox, stated precisely}
\label{sec:sandbox}

The runner is launched by one function, \code{runSandboxed}, which detects whether
\code{bwrap} (bubblewrap) is present and adapts.

On the workstation and in local development, where \code{bwrap} is installed, each
analysis runs under bubblewrap with the whole filesystem read-only
(\code{-{}-ro-bind / /}), no network (\code{-{}-unshare-net}, verified to block
outbound), a fresh \code{tmpfs} scratch, and \code{-{}-die-with-parent}, the whole
wrapped in \code{timeout}. The deployed Cloud Run image, by contrast, ships
\emph{without} \code{bwrap}. There the kernel runs in \code{timeout}-only mode --
which is exactly what \route{/health} reports as \code{sandbox:"timeout"} -- and the
isolation boundary is gVisor, the sandbox Cloud Run already places around the
container, together with the read-only image and Cloud Run's own egress posture. We
do \emph{not} claim that \code{bwrap} runs in production; it does not. On Cloud Run,
the kernel boundary is gVisor and a wall-clock timeout bounds each call, while the
no-arbitrary-code guarantee comes from Invariant~\ref{inv:no-rce}, which holds
identically in both modes.

\begin{principle}[Defence in depth, not in place of the registry]
\label{prin:depth}
The sandbox narrows what a misbehaving runner could touch; it is not the reason the
public endpoint is safe to expose. That reason is the parameterised-only registry.
Where \code{bwrap} is present we additionally pin the filesystem read-only and the
network off; where it is not, gVisor and a per-call timeout provide the kernel
boundary. The safety claim is invariant to which mode is live.
\end{principle}

Either way, one method call is bounded by \code{COMPUTE\_TIMEOUT\_S} -- 90\,s in the
deployed image, default 60\,s -- and a timeout is surfaced as an explicit error, so
a pathological input degrades to a clean \code{timeout after Ns}, never an unbounded
run.

\subsection{The JSON contract and the purity of stdout}
\label{sec:json-contract}

The data contract is deliberately minimal: one JSON object in, one JSON object out.
A request carries \code{\{method, params\}}; a runner prints exactly one JSON object
to \code{stdout} and exits. This single discipline is what lets a method run
identically across the three planes -- the cloud container, the workstation, and the
harness all spawn the same script and parse the same object. The R helpers in
\file{r/\_io.R} centralise emission through one function, \code{ce\_emit}, which
serialises the result and quits; failures emit \code{\{error: ...\}} through
\code{ce\_fail}.

The fragility this guards against is real. R routinely writes to \code{stdout} for
reasons unrelated to the result -- a warning that a KPSS p-value lies beyond the
tabulated range, a package's registration notice on load. Any such stray line would
be concatenated with the JSON and break the parse. The helpers therefore set
\code{options(warn = -1)} and suppress load messages, so that the only thing on
\code{stdout} is the result object.

Progress reporting is the interesting case, because the heavy tower estimators
genuinely need to stream progress, yet the kernel's stdout must stay pure. The
engine resolves this with a single environment gate. The helper \code{ce\_progress}
emits a newline-terminated progress line \emph{only} when the async worker has set
\code{CE\_PROGRESS=1}; under the synchronous kernel that variable is unset, so
\code{ce\_progress} is a silent no-op and the pure-single-JSON contract of
\route{/api/compute/run} is unchanged. The tower (Section~\ref{sec:tower}) sets
\code{CE\_PROGRESS=1}, reads the newline-delimited progress lines to drive its
event stream, and treats the final unterminated object as the result. The same
runner thus serves both planes without a branch in its own logic.

\subsection{The data plane}
\label{sec:data-plane}

The primary stored panel is \file{papers/contagion-channels/data/G20.xlsx}: daily
equity log-returns for 18 markets spanning 2006-01-12 to 2026-03-18, which yields
$18 \times 17 = 306$ directed market pairs for the information-flow methods. Transfer
entropy is computed on log-returns, which are $I(0)$, never on price levels, which
are $I(1)$; the stationarity gate is enforced rather than assumed. A second panel,
\code{g20\_24}, supplies the 24-series set used by the NAMH methods (19 equity
indices plus five commodity futures), read from the published \code{namh} package's
bundled data so that the cloud image and the workstation see byte-identical inputs.

Both panels are byte-hashed with SHA-256 at kernel and tower startup, and the hash
travels with the result in its provenance stamp. A reported figure therefore names
the exact data state it was computed against, which is what makes a later
reproduction meaningful: the panel is identified not by a filename but by the hash
of its bytes. Datasets that are planned but not provisioned -- an intraday G20 panel,
for instance -- are present in the registry as honest placeholders with no series
list, so \code{validate} refuses any analysis against them rather than fabricating a
result. This is the data-plane expression of the same discipline that governs the
method catalogue: a capability is exposed only when it is real.
